% SPDX-FileCopyrightText: © 2026 Idiap Research Institute <contact@idiap.ch>
%
% SPDX-FileContributor: Andre Freitas <andre.freitas@idiap.ch>
% Neuro-symbolic AI Group
%
% SPDX-FileContributor: Danilo Gusicuma <danilo.gusicuma@idiap.ch>
% Neuro-symbolic AI Group
%
% SPDX-License-Identifier: MIT
% Generated from step-explanations.json
\paragraph{Step-by-step reconstruction.}
A lower-bound knowledge base contains a highly connected complexity-class definition and a less connected proposed proof-to-circuit translation. The target requires uniformity.
The following is a source-guided illustrative reconstruction. Internal procedures are prescriptions; the native declarations do not establish their execution.
\begin{enumerate}[leftmargin=*,label=\arabic*.]
\item \textbf{Ground the route evidence} (\texttt{e0}).
\begin{description}[style=standard,labelindent=0pt,leftmargin=5.5em,labelwidth=5em,labelsep=0.5em,align=left,font=\normalfont\bfseries,itemsep=0.15em,parsep=0pt,topsep=0.3em]
\item[Purpose] Ground statements before proposing proof connections.
\item[Input] Source bundle; target; source references
\item[Method] Normalize theorem bodies and retain hypotheses, provenance and the status of each imported claim.
\item[Check] Only licensed statements may serve as established premises; topic proximity is insufficient.
\item[Output] Grounded definitions; proposed translation claims with status
\item[Next] Route formation consumes this evidence; ungrounded statements remain conjectural or excluded.
\item[Example] Record the uniform lower-bound target, the central definition, and a translation candidate whose uniformity condition has not been checked.
\end{description}
\item \textbf{Propose a useful connection} (\texttt{a0}).
\begin{description}[style=standard,labelindent=0pt,leftmargin=5.5em,labelwidth=5em,labelsep=0.5em,align=left,font=\normalfont\bfseries,itemsep=0.15em,parsep=0pt,topsep=0.3em]
\item[Purpose] Identify a target-relevant missing implication.
\item[Input] Grounded statements; dependency context
\item[Method] Compare candidate reductions and cross-component links using symbol roles and method overlap, separating proposed edges from established dependencies.
\item[Check] A route must supply proof relevance beyond graph centrality.
\item[Output] Route hypothesis: candidate proof-to-circuit translation
\item[Next] Deduction exposes its consequence; the selected bridge is still hypothetical.
\item[Example] Choose the less central translation because, if uniform, it could connect a proof lower bound to the circuit statement required by the target.
\end{description}
\item \textbf{Expose the route's consequence} (\texttt{d0}).
\begin{description}[style=standard,labelindent=0pt,leftmargin=5.5em,labelwidth=5em,labelsep=0.5em,align=left,font=\normalfont\bfseries,itemsep=0.15em,parsep=0pt,topsep=0.3em]
\item[Purpose] Expose the candidate bridge’s mathematical payoff.
\item[Input] Translation hypothesis; retained conditions
\item[Method] Derive the proposed implication while tracking the target quantifier that the translation must preserve.
\item[Check] Do not replace a uniform target with a nonuniform consequence.
\item[Output] Translation consequence, conditional on required uniformity
\item[Next] Assessment checks this consequence's relevance and missing support.
\item[Example] The route would transfer the chosen lower-bound obligation only if one uniform construction supplies the translated objects across input sizes.
\end{description}
\item \textbf{Assess the consequence} (\texttt{e1}).
\begin{description}[style=standard,labelindent=0pt,leftmargin=5.5em,labelwidth=5em,labelsep=0.5em,align=left,font=\normalfont\bfseries,itemsep=0.15em,parsep=0pt,topsep=0.3em]
\item[Purpose] Identify the decisive gap in the payoff.
\item[Input] Derived consequence; uniformity assumption
\item[Method] Compare the consequence to the original target and inspect whether graph prominence concealed a missing theorem.
\item[Check] A highly connected definition does not compensate for an unproved translation property.
\item[Output] Consequence assessment: missing uniformity-preservation support
\item[Next] Entrypoint revision consumes this diagnosis instead of repeating the full ambitious target.
\item[Example] The definition contributes vocabulary, but the route's progress depends on proving that translation does not use size-specific advice.
\end{description}
\item \textbf{Revise the entrypoint} (\texttt{a1}).
\begin{description}[style=standard,labelindent=0pt,leftmargin=5.5em,labelwidth=5em,labelsep=0.5em,align=left,font=\normalfont\bfseries,itemsep=0.15em,parsep=0pt,topsep=0.3em]
\item[Purpose] Reduce the route to a testable obligation.
\item[Input] Consequence assessment; localized uniformity gap
\item[Method] Anchor an agenda on one statement, attach its hypotheses and links, and specify a stop or pruning condition.
\item[Check] The entrypoint must directly control the missing implication.
\item[Output] Entrypoint hypothesis: uniformity-preservation lemma
\item[Next] Refutation challenges this entrypoint before a proof handoff is justified.
\item[Example] Replace the full lower-bound demand with the lemma that the proposed translation uses no nonuniform advice as input size varies.
\end{description}
\item \textbf{Challenge the entrypoint} (\texttt{r0}).
\begin{description}[style=standard,labelindent=0pt,leftmargin=5.5em,labelwidth=5em,labelsep=0.5em,align=left,font=\normalfont\bfseries,itemsep=0.15em,parsep=0pt,topsep=0.3em]
\item[Purpose] Test the selected lemma in its stated regime.
\item[Input] Uniformity-preservation entrypoint; hypotheses
\item[Method] Search for a translation instance requiring advice or a regime change; distinguish an actual counterexample from a missing argument.
\item[Check] A falsifier must attack the selected lemma, not a different complexity claim.
\item[Output] Falsification finding: advice dependence or unresolved uniformity proof
\item[Next] Final readiness joins this finding with the same entrypoint.
\item[Example] A translation that chooses unrelated advice at each input size defeats the proposed uniformity lemma, leaving any nonuniform statement separate.
\end{description}
\item \textbf{Bound the proof handoff} (\texttt{e2}).
\begin{description}[style=standard,labelindent=0pt,leftmargin=5.5em,labelwidth=5em,labelsep=0.5em,align=left,font=\normalfont\bfseries,itemsep=0.15em,parsep=0pt,topsep=0.3em]
\item[Purpose] Assess a proof task with its criticism intact.
\item[Input] Entrypoint hypothesis + its falsification assessment
\item[Method] Retain the route dossier, hub-bias diagnosis and transfer burden when judging readiness.
\item[Check] Readiness cannot arise solely from connectivity or a favorable label.
\item[Output] Readiness assessment: bounded proof task or reorientation rationale
\item[Next] Original routing needs a separate ProofQueueItem adapter; the native return does not construct it.
\item[Example] Hand off a precise unresolved uniformity obligation if viable, or stop the advice-dependent route without claiming a lower bound.
\end{description}
\end{enumerate}
\paragraph{Changed-input contrast.}
If the translation is independently shown to preserve uniformity, it can support a suitable proof entrypoint. Merely increasing the centrality of the definition never supplies that implication.
